
\documentclass[11pt]{article}
\usepackage[margin=1in]{geometry}
\usepackage{amsmath,amssymb}
\usepackage{enumitem}
\usepackage[T1]{fontenc}
\usepackage{lmodern}
\setlength{\parskip}{0.5em}
\setlength{\parindent}{0pt}

\begin{document}

\begin{center}
{\LARGE \textbf{IB Physics 2016--2020}}\\
{\large \textbf{Mapped Collection for D2 + D3}}\\
\vspace{0.3em}
\textit{Near-literal paraphrases matched to the modern electric- and magnetic-field units.}
\end{center}

\textbf{Use.} These are intentionally very close in setup, data, and subpart structure to the original old-syllabus problems, while avoiding verbatim copying.

\section*{Problems}

\begin{enumerate}[leftmargin=*, itemsep=1.1em]

\item \textbf{M16 HL Paper 2, Q6(b) --- moving charges in a conductor in a magnetic field}\\
\textit{Diagram.} A straight metallic strip has width $L$. Charge carriers move along the strip with speed $v$. A uniform magnetic field of flux density $B$ acts perpendicular to the strip, directed \emph{into} the page.
\begin{enumerate}[label=(\alph*)]
\item Show that the potential difference that appears across the width of the strip is
\[
V=vBL.
\]
\item Copy the sketch and mark clearly which side of the strip becomes negative.
\item Explain how the magnetic force on the mobile charges produces this potential difference.
\end{enumerate}

\item \textbf{M17 TZ2 HL Paper 2, Q6(a) and Q6(d) --- current in a transmission cable and force between parallel cables}\\
\textit{Diagram for (a).} The cross-section of a power cable shows many copper wires grouped together and surrounded by an insulating layer. Both regions are subject to an electric field.\\
\textit{Diagram for (d).} Two long, identical transmission cables hang parallel to one another and carry alternating current.
\begin{enumerate}[label=(\alph*)]
\item Explain why the copper part of the cable can carry a large current whereas the insulating layer does not.
\item In terms of charge carriers, compare what happens inside the copper with what happens inside the insulator.
\item During one complete ac cycle, describe how the magnetic force between the two parallel cables varies.
\item State when the force is zero, when its magnitude is greatest, and whether the cables attract or repel at those times.
\end{enumerate}

\item \textbf{M18 TZ2 HL Paper 2, Q8 --- thundercloud and Earth modelled as a capacitor}\\
\textit{Diagram.} The negatively charged base of a thundercloud is shown above a flat region of Earth. The cloud base and the Earth's surface are modelled as the two plates of a parallel-plate capacitor. The area of the cloud base is $1.2\times10^{8}\,\text{m}^2$, the separation is $1600\,\text{m}$, the charge on the cloud is $-25\,\text{C}$, and the permittivity of air is taken as $8.8\times10^{-12}\,\text{F m}^{-1}$.
\begin{enumerate}[label=(\alph*)]
\item Show that the capacitance of this cloud-Earth system is about $6.6\times10^{-7}\,\text{F}$.
\item Find the potential difference between the thundercloud and Earth.
\item Determine the energy stored in the electric field.
\item During a lightning discharge, take the time constant of the discharge to be $32\,\text{ms}$ and the duration of the strike to be $18\,\text{ms}$. Show that the charge transferred is about $11\,\text{C}$.
\item Hence estimate the average current during the strike.
\item State one assumption needed to justify modelling the cloud-Earth system as a parallel-plate capacitor.
\end{enumerate}

\item \textbf{M19 TZ2 HL Paper 2, Q5 --- proton in uniform magnetic field}\\
\textit{Diagram.} A proton travels in a circular path inside a region where the magnetic field is uniform and directed into the plane of the page.
\begin{enumerate}[label=(\alph*)]
\item Add an arrow for the proton's instantaneous velocity and an arrow for the magnetic force at a chosen point on the orbit.
\item The proton speed is $2.16\times10^{6}\,\text{m s}^{-1}$ and the magnetic flux density is $0.042\,\text{T}$. Calculate the radius of the orbit.
\item Explain why the path is circular even though the proton speed stays constant.
\end{enumerate}

\item \textbf{M19 TZ2 HL Paper 2, Q9(b) --- field lines and the zero-field point}\\
\textit{Diagram.} Two fixed charges, $X$ and $Y$, lie on a straight line. A point $P$ lies on the same line, outside the interval between the charges. A field-line sketch is given for the pair. The separation $XY$ is $0.600\,\text{m}$ and the distance $PY$ is $0.820\,\text{m}$. The electric field at $P$ is zero.
\begin{enumerate}[label=(\alph*)]
\item Use the zero-field condition to determine, to one significant figure, the ratio $Q/q$, where $Q$ and $q$ are the magnitudes of the charges on $X$ and $Y$ respectively.
\item From the field-line pattern, infer whether the two charges have the same sign or opposite signs, and justify your answer briefly.
\end{enumerate}

\item \textbf{N19 HL Paper 2, Q9 --- two capacitors sharing charge}\\
\textit{Diagram.} Capacitor $X$ has capacitance $18\,\mu\text{F}$. Initially, one plate of $X$ carries charge $48\,\mu\text{C}$. Capacitor $Y$ has capacitance $12\,\mu\text{F}$ and is initially uncharged. The capacitors are then connected through a resistor by closing a switch, so charge redistributes until both capacitors have the same potential difference.
\begin{enumerate}[label=(\alph*)]
\item With the switch open, calculate the energy stored in capacitor $X$.
\item After the switch is closed, determine the final charge on capacitor $X$ and the final charge on capacitor $Y$.
\item Find the total electrical energy stored after equilibrium is reached.
\item Explain why the final stored energy is less than the initial stored energy.
\end{enumerate}

\item \textbf{N20 HL Paper 2, Q8 --- electric potential around a positively charged sphere}\\
\textit{Diagram.} A positively charged conducting sphere is shown with points $A$ and $B$ at different distances from its centre, and a graph is to be sketched against radial distance $r$.
\begin{enumerate}[label=(\alph*)]
\item Compare the work done per unit charge in bringing a positive test charge from infinity to point $A$ and from infinity to point $B$.
\item Sketch the variation of electric potential $V$ with distance $r$ from the centre of the sphere, showing the behaviour inside the sphere and outside it.
\item A particle has electric potential energy $1.7\times10^{-16}\,\text{J}$ when it is at a point where the electric potential is $1.1\times10^{3}\,\text{V}$. Determine the charge of the particle.
\item Use the electric-potential expression for a charged sphere or point charge, together with the given distance information, to determine the charge on the sphere.
\item State one reason why it is useful to compare electric fields with gravitational fields.
\end{enumerate}

\end{enumerate}

\end{document}
